% GENERATED FILE. Edit canonical lessons, then run npm run generate:books.
% canonical-source-hash: fnv1a64:02aa904b29c6a6bb
% canonical-lessons: MW-C111-chavno, MW-C111-ichchha, MW-C111-man, MW-C111-as, MW-C111-mag

\chapter{To want, Wish, Mind, Hope, Demand}
\label{ch:mw-to-want-wish-mind-hope-demand}
\hlchaptermodality{111}

\begin{chapteropening}
\textbf{By the end of this chapter:} \emph{I can say to want, a wish, the mind, hope and a demand.}

Complete the last lesson of chapter 111: I can say to want, a wish, the mind, hope and a demand.
\end{chapteropening}

\section[\texorpdfstring{\mw{चावणो} (chāvṇo)}{चावणो (chāvṇo)}]{\mw{चावणो} (chāvṇo) — to want}

\label{lesson:MW-C111-chavno}



\begin{quote}
Before the new one: say the Marwadi for to fly, then the Marwadi for to weigh.
\end{quote}

\subsection*{You'll want to know: \mw{चावणो}}
\textbf{\mw{चावणो}} — \emph{chāvṇo} — \textquotedblleft{}to want\textquotedblright{}.

\begin{cousinweb}
\textbf{\mw{चावणो}} (\emph{chāvṇo}) is a tadbhava, a word inherited and reshaped by sound change. It comes from \emph{*cāh}, \textquotedblleft{}see, look for, desire\textquotedblright{}, a form CDIAL reconstructs from its descendants. CDIAL lists Old Marwari \emph{cāhaï} under it (CDIAL 4775). CDIAL calls the further source of this reconstructed verb uncertain.
\end{cousinweb}

\begin{grammarlens}[title={What you've built}]
One more word for this chapter.
\end{grammarlens}

\subsection*{Your turn}
\begin{itemize}
\raggedright
  \item \emph{Say it:} \emph{chāvṇo}
  \item \emph{Say it:} \emph{chāvṇo}, once more
  \item \emph{Say it:} say \emph{chāvṇo}
  \item \emph{Recall:} say the Marwadi for to fly, then the Marwadi for to weigh, then say \emph{chāvṇo} again
\end{itemize}

\begin{tcolorbox}[breakable,colback=teal!4,colframe=teal!35!black,title={Before you move on}]
What does \mw{चावणो} mean? (\textquotedblleft{}To want\textquotedblright{}.) Say it once more.
\end{tcolorbox}
\section[\texorpdfstring{\mw{इच्छा} (ichchhā)}{इच्छा (ichchhā)}]{\mw{इच्छा} (ichchhā) — a wish}

\label{lesson:MW-C111-ichchha}



\begin{quote}
Before the new one: say the Marwadi for to weigh, then the Marwadi for to want.
\end{quote}

\subsection*{You'll want to know: \mw{इच्छा}}
\textbf{\mw{इच्छा}} — \emph{ichchhā} — \textquotedblleft{}a wish\textquotedblright{}.

\begin{cousinweb}
\textbf{\mw{इच्छा}} (\emph{ichchhā}) is Sanskrit \emph{icchā}, \textquotedblleft{}wish\textquotedblright{}, taken over whole as a learned loan, a tatsama (CDIAL 1557).
\end{cousinweb}

\begin{grammarlens}[title={What you've built}]
One more word for this chapter.
\end{grammarlens}

\subsection*{Your turn}
\begin{itemize}
\raggedright
  \item \emph{Say it:} \emph{ichchhā}
  \item \emph{Say it:} \emph{ichchhā}, once more
  \item \emph{Say it:} say \emph{ichchhā}
  \item \emph{Recall:} say the Marwadi for to weigh, then the Marwadi for to want, then say \emph{ichchhā} again
\end{itemize}

\begin{tcolorbox}[breakable,colback=teal!4,colframe=teal!35!black,title={Before you move on}]
What does \mw{इच्छा} mean? (\textquotedblleft{}A wish\textquotedblright{}.) Say it once more.
\end{tcolorbox}
\section[\texorpdfstring{\mw{मन} (man)}{मन (man)}]{\mw{मन} (man) — the mind, a wish}

\label{lesson:MW-C111-man}



\begin{quote}
Before the new one: say the Marwadi for to want, then the Marwadi for a wish.
\end{quote}

\subsection*{You'll want to know: \mw{मन}}
\textbf{\mw{मन}} — \emph{man} — \textquotedblleft{}the mind, a wish\textquotedblright{}.

\begin{cousinweb}
\textbf{\mw{मन}} (\emph{man}) is Sanskrit \emph{manas}, \textquotedblleft{}mind\textquotedblright{}, taken over whole as a learned loan, a tatsama (CDIAL 9822). The \emph{-as} ending was dropped. English \emph{mind} is a distant cousin.
\end{cousinweb}

\begin{grammarlens}[title={What you've built}]
One more word for this chapter.
\end{grammarlens}

\subsection*{Your turn}
\begin{itemize}
\raggedright
  \item \emph{Say it:} \emph{man}
  \item \emph{Say it:} \emph{man}, once more
  \item \emph{Say it:} say \emph{man}
  \item \emph{Recall:} say the Marwadi for to want, then the Marwadi for a wish, then say \emph{man} again
\end{itemize}

\begin{tcolorbox}[breakable,colback=teal!4,colframe=teal!35!black,title={Before you move on}]
What does \mw{मन} mean? (\textquotedblleft{}The mind, a wish\textquotedblright{}.) Say it once more.
\end{tcolorbox}
\section[\texorpdfstring{\mw{आस} (ās)}{आस (ās)}]{\mw{आस} (ās) — hope}

\label{lesson:MW-C111-as}



\begin{quote}
Before the new one: say the Marwadi for a wish, then the Marwadi for the mind.
\end{quote}

\subsection*{You'll want to know: \mw{आस}}
\textbf{\mw{आस}} — \emph{ās} — \textquotedblleft{}hope\textquotedblright{}.

\begin{cousinweb}
\textbf{\mw{आस}} (\emph{ās}) is a tadbhava, a word inherited from Sanskrit and reshaped by sound change. It comes from \emph{āśā}, \textquotedblleft{}wish, expectation\textquotedblright{}. CDIAL lists Old Marwari \emph{āsa} under it (CDIAL 1456).
\end{cousinweb}

\begin{grammarlens}[title={What you've built}]
One more word for this chapter.
\end{grammarlens}

\subsection*{Your turn}
\begin{itemize}
\raggedright
  \item \emph{Say it:} \emph{ās}
  \item \emph{Say it:} \emph{ās}, once more
  \item \emph{Say it:} say \emph{ās}
  \item \emph{Recall:} say the Marwadi for a wish, then the Marwadi for the mind, then say \emph{ās} again
\end{itemize}

\begin{tcolorbox}[breakable,colback=teal!4,colframe=teal!35!black,title={Before you move on}]
What does \mw{आस} mean? (\textquotedblleft{}Hope\textquotedblright{}.) Say it once more.
\end{tcolorbox}
\section[\texorpdfstring{\mw{मांग} (mā̃g)}{मांग (mā̃g)}]{\mw{मांग} (mā̃g) — a demand}

\label{lesson:MW-C111-mag}



\begin{quote}
Before the new one: say the Marwadi for the mind, then the Marwadi for hope.
\end{quote}

\subsection*{You'll want to know: \mw{मांग}}
\textbf{\mw{मांग}} — \emph{mā̃g} — \textquotedblleft{}a demand\textquotedblright{}.

\begin{cousinweb}
\textbf{\mw{मांग}} (\emph{mā̃g}) is a tadbhava, a word inherited from Sanskrit and reshaped by sound change. It comes from \emph{mārgya}, \textquotedblleft{}to be sought\textquotedblright{}. CDIAL lists Old Marwari \emph{māṃga} under it (CDIAL 10078). \textbf{\mw{मांग}} (\emph{mā̃g}) is the noun of \emph{māṅgṇo}, \textquotedblleft{}to ask for\textquotedblright{}; the Old Marwari row means \textquotedblleft{}betrothed\textquotedblright{}.
\end{cousinweb}

\begin{grammarlens}[title={What you've built}]
One more word for this chapter.
\end{grammarlens}

\subsection*{Your turn}
\begin{itemize}
\raggedright
  \item \emph{Say it:} \emph{mā̃g}
  \item \emph{Say it:} \emph{mā̃g}, once more
  \item \emph{Say it:} say \emph{mā̃g}
  \item \emph{Recall:} say the Marwadi for the mind, then the Marwadi for hope, then say \emph{mā̃g} again
\end{itemize}

\begin{tcolorbox}[breakable,colback=teal!4,colframe=teal!35!black,title={Before you move on}]
What does \mw{मांग} mean? (\textquotedblleft{}A demand\textquotedblright{}.) Say it once more.
\end{tcolorbox}
